\chapter{ТЕОРЕТИЧЕСКИЕ ОСНОВЫ МОДЕЛИРОВАНИЯ ТЕРМОУПРУГИХ ВОЛН В ГЕОЛОГИЧЕСКИХ СРЕДАХ}

\section{Физика термоупругих волн в геологических средах}

Термоупругие волны представляют собой волновые процессы в деформируемой среде, при которых механические смещения и температурные возмущения оказываются взаимосвязанными. В отличие от чисто упругой постановки, где динамика среды определяется только законами сохранения импульса и упругими соотношениями между напряжениями и деформациями, термоупругая модель учитывает двустороннюю связь между тепловым и механическим полями. С одной стороны, изменение температуры вызывает тепловое расширение или сжатие материала, вследствие чего возникают дополнительные напряжения. С другой стороны, изменение объёма и деформационного состояния среды влияет на распределение температуры через термомеханическую генерацию и перераспределение тепла. Именно такая взаимосвязь особенно существенна для геологических сред, где неоднородность строения, изменчивость параметров, наличие порового пространства, фазовых границ и температурных градиентов приводят к сложной многополевой динамике.

В качестве базовой математической постановки термоупругой задачи используется система связанных уравнений Куна--Томпсона. Базовое уравнение движения имеет вид
\begin{equation}
\rho \frac{\partial^2 u_i}{\partial t^2}=\frac{\partial \sigma_{ij}}{\partial x_j}.
\end{equation}

Данное соотношение выражает локальный баланс линейного импульса в сплошной среде. Здесь $u_i$~--- компоненты вектора смещения, то есть функции, описывающие отклонение материальной точки от исходного положения; $\rho$~--- плотность среды; $\sigma_{ij}$~--- компоненты тензора напряжений; индекс $i$ указывает направление соответствующей компоненты смещения, а индекс $j$ связан с пространственной производной по координате $x_j$. Левая часть уравнения характеризует инерционную реакцию среды: произведение плотности на ускорение материальной точки. Правая часть описывает внутренние силы, возникающие из-за неоднородного распределения напряжений. Следовательно, данное уравнение связывает кинематику смещений с внутренним напряжённым состоянием материала.

Однако само по себе уравнение движения не замыкает задачу, поскольку тензор напряжений $\sigma_{ij}$ должен быть выражен через деформации и температуру. Для этого используется термоупругое конститутивное соотношение
\begin{equation}
\sigma_{ij}=\lambda \delta_{ij}\varepsilon_{kk}+2\mu \varepsilon_{ij}-\gamma \delta_{ij}T.
\end{equation}

Здесь $\varepsilon_{ij}$~--- компоненты тензора малых деформаций, $\delta_{ij}$~--- символ Кронекера, $\lambda$ и $\mu$~--- параметры Ламе, $T$~--- изменение температуры, $\gamma$~--- коэффициент термоупругой связи. Физический смысл этого соотношения состоит в том, что полное напряжение складывается из трёх частей. Первый член, $\lambda \delta_{ij}\varepsilon_{kk}$, описывает вклад объёмной деформации: след тензора деформаций $\varepsilon_{kk}$ соответствует относительному изменению объёма, а параметр $\lambda$ характеризует чувствительность среды к объёмному сжатию или расширению. Второй член, $2\mu\varepsilon_{ij}$, отвечает за сдвиговые деформации; параметр $\mu$ является модулем сдвига и определяет сопротивление среды сдвиговым искажениям формы. Третий член, $-\gamma \delta_{ij}T$, непосредственно отражает тепловую природу термоупругости: повышение температуры вызывает стремление материала к свободному расширению, а при наличии ограничений это расширение преобразуется в дополнительные напряжения.

Для корректной интерпретации выражения необходимо пояснить и тензор деформаций. В предположении малых деформаций он записывается как
\begin{equation}
\varepsilon_{ij}=\frac{1}{2}\left(\frac{\partial u_i}{\partial x_j}+\frac{\partial u_j}{\partial x_i}\right).
\end{equation}

Такое определение означает, что модель строится в предположении малых перемещений и малых поворотов, когда геометрическая нелинейность может быть опущена. Это допущение является типичным для задач линейной волновой динамики, если рассматриваемые возмущения не приводят к крупным разрушениям и пластическим течениям.

Третьим элементом системы служит уравнение теплопереноса с учётом термомеханической связи:
\begin{equation}
k\nabla^2T-C\frac{\partial T}{\partial t}=\gamma T_0\frac{\partial \varepsilon_{kk}}{\partial t}.
\end{equation}

В этом уравнении $k$~--- коэффициент теплопроводности, $C$~--- теплоёмкость, $T_0$~--- средняя или опорная температура, $\nabla^2$~--- оператор Лапласа, $\varepsilon_{kk}$~--- объёмная деформация. Левая часть описывает тепловую эволюцию: член $k\nabla^2T$ характеризует пространственное выравнивание температуры за счёт теплопроводности, а член $-C\partial T/\partial t$ отражает инерционность теплового режима, то есть необходимость затраты энергии на изменение температуры среды. Правая часть, $\gamma T_0 \, \partial \varepsilon_{kk}/\partial t$, отвечает за связь с механическим полем: изменение объёма во времени сопровождается термомеханическим источником тепла или его поглощением. Следовательно, уравнение показывает, что температура в термоупругой среде меняется не только из-за теплопроводности, но и вследствие динамики деформирования.

Таким образом, уравнения (1.1)--(1.4) образуют единую систему связанных дифференциальных уравнений в частных производных. Уравнение движения определяет, как поле напряжений порождает ускорение и смещения. Конститутивное соотношение выражает напряжение через деформацию и температуру. Уравнение теплопереноса замыкает систему, связывая тепловое поле с объёмной деформацией. Именно такая структура отражает физическую сущность термоупругих волн: механическое возмущение порождает температурный отклик, а температурный отклик, в свою очередь, меняет напряжённое состояние и характер распространения механической волны.

В неоднородной геологической среде параметры $\rho$, $\lambda$, $\mu$, $\gamma$ и, вообще говоря, $k$ и $C$, перестают быть постоянными и зависят от координат. В этом случае уравнение движения принимает вид
\begin{equation}
\rho(x)\frac{\partial^2 u_i}{\partial t^2}
=
\frac{\partial}{\partial x_j}
\left[
\lambda(x)\delta_{ij}\frac{\partial u_k}{\partial x_k}
+
\mu(x)\left(\frac{\partial u_i}{\partial x_j}+\frac{\partial u_j}{\partial x_i}\right)
\right]
-
\gamma(x)\frac{\partial T}{\partial x_i}.
\end{equation}

Это уравнение принципиально важнее для геофизических приложений, чем однородная модель. Его левая часть по-прежнему описывает инерцию, но теперь она зависит от пространственно изменяющейся плотности $\rho(x)$. В правой части члены с $\lambda(x)$ и $\mu(x)$ отражают неоднородность упругих свойств: волна в разных точках пространства распространяется в материале с различной жёсткостью. Последний член $-\gamma(x)\partial T/\partial x_i$ показывает, что температурный градиент становится дополнительным источником механического воздействия. Иными словами, не только температура как таковая, но и её пространственная неоднородность изменяет динамическое состояние среды.

Если среда анизотропна, скалярных коэффициентов Ламе уже недостаточно. Тогда используется тензор жёсткости:
\begin{equation}
\sigma_{ij}=C_{ijkl}\varepsilon_{kl}-\gamma_{ij}T.
\end{equation}

Здесь $C_{ijkl}$~--- тензор четвёртого ранга, задающий анизотропные упругие свойства; $\gamma_{ij}$~--- тензор термоупругой связи. В изотропной среде тензор $C_{ijkl}$ сводится к комбинации параметров Ламе, однако для реальных пород, особенно слоистых и текстурированных, жёсткость существенно зависит от направления. Это означает, что скорость, затухание и поляризация волн становятся ориентационно зависимыми. Следовательно, анизотропия усложняет не только физическую интерпретацию, но и численную реализацию модели, поскольку каждое направление распространения волны может демонстрировать различную динамику даже в одной и той же точке среды.

Для интерпретации механической части модели полезно сопоставить её с классическими уравнениями упругой динамики. Для изотропного случая и при отсутствии объёмных сил уравнение движения записывается в векторной форме как
\begin{equation}
\rho\,\ddot{\mathbf{u}}=(\lambda+\mu)\nabla(\nabla\!\cdot\mathbf{u})+\mu\nabla^2\mathbf{u},
\end{equation}
или эквивалентно
\begin{equation}
\rho\,\ddot{\mathbf{u}}=(\lambda+2\mu)\nabla(\nabla\!\cdot\mathbf{u})-\mu\nabla\times(\nabla\times\mathbf{u}).
\end{equation}

Эти записи показывают, что динамика упругого поля естественным образом распадается на продольную и поперечную составляющие. Продольная часть связана с дивергенцией $\nabla\!\cdot\mathbf{u}$, то есть с локальным изменением объёма, а поперечная~--- с вихревой частью $\nabla\times\mathbf{u}$, то есть с сдвиговыми движениями без изменения объёма. Для термоупругой модели это особенно важно, поскольку температурная связь непосредственно входит именно в объёмную часть деформирования через $\varepsilon_{kk}$. Следовательно, тепловое поле сильнее всего взаимодействует с продольными волнами, тогда как поперечные волны вовлекаются в термоупругий процесс главным образом косвенно~--- через неоднородность среды, граничные эффекты и связь мод на интерфейсах.

Из классической упругой модели следуют скорости продольных и поперечных волн:
\begin{equation}
c_P=\sqrt{\frac{\lambda+2\mu}{\rho}},\qquad
c_S=\sqrt{\frac{\mu}{\rho}}.
\end{equation}

Из этих выражений видно, что скорость продольных волн определяется как модулем сдвига, так и объёмной жёсткостью, тогда как скорость поперечных волн~--- только модулем сдвига. Это имеет прямой геофизический смысл. Параметры $\lambda$ и $\mu$, а также плотность $\rho$, зависят от минерального состава, пористости, трещиноватости, насыщенности и температурного состояния породы. Поэтому пространственные вариации этих характеристик вызывают изменение скоростей волн, перераспределение энергии, рефракцию, частичное отражение и рассеяние.

Математическая постановка термоупругой задачи не ограничивается самой системой PDE; она должна быть дополнена начальными и граничными условиями. Начальные условия определяют состояние среды в начальный момент времени:
\begin{equation}
u_i(x,0)=u_i^{(0)}(x),\qquad
\frac{\partial u_i}{\partial t}(x,0)=v_i^{(0)}(x),
\end{equation}
а также
\begin{equation}
T(x,0)=T^{(0)}(x).
\end{equation}

Если используется тепловая модель с первым порядком по времени, этих условий достаточно. Физически они задают исходное деформированное и тепловое состояние среды, из которого развивается волновой процесс.

Граничные условия в термоупругой задаче имеют двойственную природу. Для механической части можно задавать либо смещения, либо поверхностные силы. Для свободной поверхности справедливо условие отсутствия тракции
\begin{equation}
\boldsymbol{\sigma}\cdot \mathbf{n}=\mathbf{0},
\end{equation}
где $\mathbf{n}$~--- внешняя нормаль к границе. Это условие моделирует, например, поверхность геологического массива, не нагруженную внешним давлением. На контакте двух сред обычно требуется непрерывность смещений и непрерывность нормальных и касательных компонент напряжения, если предполагается идеальное сцепление. Для тепловой части аналогично можно задавать либо саму температуру на границе, либо тепловой поток:
\begin{equation}
T|_{\partial\Omega}=T_b,\qquad
-k\nabla T\cdot \mathbf{n}=q_b.
\end{equation}

Первое условие соответствует фиксированной температуре, второе~--- заданному притоку или оттоку тепла. Различие между механическими и тепловыми краевыми условиями принципиально: механические условия управляют отражением и передачей волновой энергии на границе, тогда как тепловые условия влияют на степень диссипации и характер температурной релаксации.

Выбор граничных условий существенно меняет решение. Свободная поверхность вызывает преобразование типов волн и образование отражённых компонент. Жёстко закреплённая граница усиливает локализацию напряжений. Контраст температурных условий на границе меняет пространственную структуру теплового поля и тем самым~--- термоупругую добавку к напряжению. В неоднородной среде это приводит к особенно сложной картине, поскольку волна многократно взаимодействует с интерфейсами, включениями и зонами контрастных свойств.

Для геологических сред характерны отражение, преломление, рассеяние и частичная локализация энергии. На границе двух слоёв с разными импедансами $Z=\rho c$ возникает частичное отражение и частичная передача волны. Если граница имеет криволинейную форму, а параметры меняются плавно, наблюдается рефракция: траектория переноса энергии изгибается в сторону области меньшей или большей волновой скорости. Если же неоднородности имеют размеры, сравнимые с длиной волны, усиливается рассеяние. В термоупругой постановке к этому добавляется ещё и затухание, связанное с перераспределением части механической энергии в тепловую. Поэтому термоупругая волна в реальной породе~--- это не просто упругая волна с небольшой поправкой, а сложный многополевой процесс, в котором динамика смещений, напряжений и температуры должна рассматриваться совместно.

Итак, физика термоупругих волн в геологических средах определяется тремя ключевыми обстоятельствами. Во-первых, механическое и тепловое поля не являются независимыми. Во-вторых, геологическая среда пространственно неоднородна и часто анизотропна. В-третьих, волновой процесс всегда зависит не только от внутренних коэффициентов модели, но и от начальных и краевых условий. Следовательно, уже на уровне постановки задачи становится очевидно, что даже линейная термоупругая модель приводит к вычислительно сложной и содержательно богатой системе PDE.

\section{Геологические среды и их свойства}

Геологическая среда не может рассматриваться как простой однородный континуум с постоянными коэффициентами. Даже в пределах одного литологического типа свойства породы могут существенно изменяться в зависимости от минералогического состава, пористости, трещиноватости, степени водонасыщения, глубины залегания, температурного режима и предыстории формирования. Поэтому при моделировании термоупругих волн геологический массив выступает как многопараметрический объект, а его свойства непосредственно определяют коэффициенты математической модели.

Наиболее простой идеализацией является однородная изотропная среда. В этом случае параметры $\rho$, $\lambda$, $\mu$, $k$, $C$, $\gamma$ считаются постоянными, а волновые фронты в отсутствии сложной геометрии распространяются предсказуемо. Такая модель удобна для аналитического анализа, вывода дисперсионных соотношений, тестирования численных схем и формирования базовых обучающих выборок. Однако для реальных геологических приложений она, как правило, недостаточна.

Более реалистичной является неоднородная среда, в которой свойства зависят от координат. Неоднородность может быть слоистой, когда параметры меняются преимущественно вдоль одной координаты; локализованной, когда в основной массе породы присутствуют включения или линзы с иными характеристиками; либо мультиструктурной, когда среда имеет иерархию неоднородностей разных масштабов. Именно здесь уравнение (1.5) приобретает принципиальное значение, поскольку переменные коэффициенты входят в дифференциальный оператор и изменяют характер решения уже на уровне самой структуры PDE.

Особый случай составляет анизотропная среда. В ней механический отклик зависит от направления деформирования. Типичным примером являются сланцы, гнейсы, слоистые осадочные толщи и породы с выраженной системой трещин. В такой ситуации тензор жёсткости $C_{ijkl}$ уже нельзя свести к двум параметрам Ламе, а термоупругая связь также может быть тензорной, как в формуле (1.6). Это означает, что одинаковые по амплитуде воздействия, приложенные в разных направлениях, приводят к различным напряжённо-деформированным и тепловым откликам. Для волновой динамики отсюда следует различие фазовых скоростей, направленная зависимость поляризации и возможность расщепления мод.

Плотность $\rho$ входит в уравнение движения как инерционный коэффициент. Чем выше плотность, тем больше сопротивление среды ускорению при прочих равных условиях. Но влияние плотности не изолировано: она входит и в выражения для волновых скоростей через отношения жёсткости к массе. Следовательно, плотность влияет не только на амплитудный отклик, но и на кинематику распространения волн. В геологических средах плотность обычно коррелирует с пористостью, насыщенностью и минеральным составом, поэтому её пространственные вариации часто служат индикатором глубинных структур и границ слоёв.

Упругие свойства в изотропной модели задаются параметрами Ламе $\lambda$ и $\mu$. Модуль сдвига $\mu$ контролирует способность породы сопротивляться изменению формы. Если $\mu$ мал, то поперечные волны замедляются, а при стремлении $\mu$ к нулю среда теряет способность поддерживать сдвиговые напряжения. Параметр $\lambda$ связан с объёмным откликом и вместе с $\mu$ определяет скорость продольных волн. Следовательно, контрасты $\lambda$ и $\mu$ между слоями вызывают различную чувствительность для $P$- и $S$-компонент и влияют на коэффициенты отражения, преломления и конверсии волн.

Тепловые свойства входят в модель через коэффициент теплопроводности $k$, теплоёмкость $C$ и коэффициент температурного расширения, который в термоупругой записи проявляется через коэффициент связи $\gamma$. Коэффициент теплопроводности определяет, насколько быстро температурный градиент сглаживается в пространстве. Чем больше $k$, тем эффективнее перераспределение тепла и тем слабее локальные температурные концентрации при прочих равных условиях. Теплоёмкость $C$ характеризует способность среды аккумулировать тепловую энергию: при большом $C$ одинаковый тепловой приток вызывает меньший рост температуры. Коэффициент теплового расширения отражает степень изменения линейных размеров материала при нагревании; именно он связывает температурное поле с механической деформацией.

По табличным данным коэффициентов линейного теплового расширения видно, что даже для пород и минеральных материалов одного класса величины могут заметно отличаться. Так, для гранита коэффициент составляет порядка $7.9\cdot 10^{-6}\ ^\circ\mathrm{C}^{-1}$, для известняка~--- около $8\cdot 10^{-6}\ ^\circ\mathrm{C}^{-1}$, для песчаника~--- около $11.6\cdot 10^{-6}\ ^\circ\mathrm{C}^{-1}$, для мрамора приведён диапазон $5.5$--$14.1\cdot 10^{-6}\ ^\circ\mathrm{C}^{-1}$, а для кварца~--- существенно меньшие значения порядка $0.77$--$1.4\cdot 10^{-6}\ ^\circ\mathrm{C}^{-1}$. Эти различия имеют прямой физический смысл. Чем выше коэффициент теплового расширения, тем сильнее механический отклик на одно и то же температурное возмущение. Следовательно, в породах с большим температурным расширением локальный нагрев способен вызывать более интенсивные термоупругие напряжения, особенно если деформация ограничена соседними участками массива.

Не менее показательны и табличные данные по теплопроводности и теплоёмкости пород. Из них видно, что, например, для кварцита теплопроводность может достигать примерно $2.6$--$7.6\ \mathrm{Вт/(м\cdot K)}$, для гранита~--- порядка $1.1$--$3.9\ \mathrm{Вт/(м\cdot K)}$, для базальта~--- около $0.4$--$3.5\ \mathrm{Вт/(м\cdot K)}$, для известняка в зависимости от состояния~--- примерно $0.7$--$4.4\ \mathrm{Вт/(м\cdot K)}$. Там же приведены интервалы удельной теплоёмкости и температуропроводности, что дополнительно подчёркивает сильную вариативность теплового поведения пород. Следовательно, даже при одинаковой геометрии и близких упругих свойствах два участка массива могут демонстрировать различный темп тепловой релаксации, а значит~--- различную степень термоупругого затухания и различную продолжительность связанного механотеплового отклика. По этой причине коэффициенты $k$ и $C$ нельзя трактовать как второстепенные: они принципиально влияют на динамику уравнения (1.4).

Существенную роль играет пористость. Для среднезернистого песчаника общий диапазон пористости составляет примерно $0.14$--$0.49$ при среднем значении около $0.34$, а эффективная пористость~--- $0.12$--$0.41$ при среднем значении около $0.27$. Для известняка общий диапазон составляет примерно $0.07$--$0.56$ при среднем значении около $0.30$, а эффективная пористость~--- от значений, близких к нулю, до $0.36$ при среднем около $0.14$. Для базальта приведён диапазон общей пористости $0.03$--$0.35$ со средним около $0.17$. Пористость влияет на свойства двояко. Во-первых, она уменьшает эффективную жёсткость и часто снижает плотность скелета, тем самым влияя на скорости волн. Во-вторых, она изменяет тепловые свойства, особенно при наличии флюидонасыщения. Поры, заполненные водой, газом или нефтью, создают дополнительную внутреннюю неоднородность, что усиливает рассеяние и приводит к локальной концентрации напряжений и температурных градиентов.

Связь между физическими свойствами среды и коэффициентами модели можно проследить непосредственно. Плотность $\rho$ и упругие коэффициенты $\lambda,\mu$ входят в динамическую часть уравнения и определяют инерцию и волновые скорости. Коэффициенты $k$ и $C$ входят в тепловой оператор и задают скорость диффузии температуры. Коэффициент температурного расширения через $\gamma$ управляет интенсивностью связи между тепловым и механическим полями. Пористость, трещиноватость и насыщенность не всегда входят в уравнение явно, но влияют на эффективные значения всех перечисленных коэффициентов. Следовательно, геологические свойства породы~--- это не внешний комментарий к модели, а её параметрическая сущность.

Контрасты свойств между слоями усложняют задачу и физически, и вычислительно. Если соседние слои резко различаются по плотности и жёсткости, то возникают сильные отражения и преобразования волн на интерфейсах. Если при этом различаются теплопроводность и коэффициент теплового расширения, то механический и тепловой отклик на границе также становятся неравномерными. В слоистых средах даже простое возмущение может порождать многократные отражения, интерференцию и запаздывающий тепловой отклик. В анизотропных средах к этому добавляется зависимость отклика от направления. Поэтому аналитическое решение подобных задач возможно лишь в редких идеализированных случаях, а численное моделирование требует либо очень тонкой сетки, либо специальных вычислительных подходов.

Таким образом, свойства геологических сред непосредственно определяют постановку задачи моделирования термоупругих волн. Они входят в коэффициенты математической модели, изменяют тип решения, влияют на устойчивость и точность численных методов и определяют общую вычислительную сложность. Чем сильнее неоднородность, слоистость и анизотропия, тем менее применимы простые аналитические модели и тем более востребованы продвинутые численные и интеллектуальные методы.

\section{Классические методы моделирования}

Классические методы численного моделирования составляют фундамент вычислительной геофизики. Именно они позволяют получить эталонные решения системы уравнений термоупругости, исследовать влияние параметров среды и формировать синтетические данные для последующего обучения моделей искусственного интеллекта. В задачах термоупругих волн основная трудность состоит в том, что приходится одновременно учитывать гиперболическую динамику механических смещений и тепловую эволюцию, имеющую выраженный диффузионный характер. Кроме того, геологическая среда, как правило, неоднородна, а её границы могут иметь сложную форму. Поэтому выбор численного метода определяется не только видом уравнений, но и геометрией области, контрастами коэффициентов и требуемой точностью.

\subsection*{Метод конечных разностей во временной области (FDTD)}

Идея FDTD состоит в замене непрерывных производных конечными разностями на дискретной сетке по пространству и времени. Пусть для простоты рассматривается одномерная переменная $u(x,t)$, заданная в узлах $x_i=i\Delta x$, $t^n=n\Delta t$. Тогда вторая производная по времени аппроксимируется как
\begin{equation}
\frac{\partial^2 u}{\partial t^2}\Big|_{i}^{n}
\approx
\frac{u_i^{n+1}-2u_i^n+u_i^{n-1}}{\Delta t^2},
\end{equation}
а вторая производная по пространству~---
\begin{equation}
\frac{\partial^2 u}{\partial x^2}\Big|_{i}^{n}
\approx
\frac{u_{i+1}^n-2u_i^n+u_{i-1}^n}{\Delta x^2}.
\end{equation}

Физический смысл такой дискретизации заключается в том, что непрерывная среда заменяется сеточным аналогом, а эволюция поля вычисляется пошагово во времени. Для системы термоупругости аналогичным образом дискретизируются уравнение движения и уравнение теплопереноса. Вектор смещений, тензорные деформации и температура обновляются на каждом временном шаге по явным или неявным формулам.

Преимущество FDTD состоит в прозрачности реализации и естественности для волновых задач. Метод непосредственно воспроизводит динамику фронтов, отражений и интерференции. Для прямоугольных областей и регулярных сеток он особенно удобен. В случае изотропной среды и простых граничных условий конечные разности позволяют эффективно реализовать явные схемы, где новое состояние системы вычисляется по известным значениям на предыдущих слоях.

Однако переход от непрерывной постановки к дискретной требует соблюдения условий устойчивости. Классическим критерием служит условие CFL, согласно которому шаг по времени должен быть достаточно мал по отношению к шагу сетки и максимальной скорости распространения возмущений:
\begin{equation}
\Delta t \le C_{\mathrm{CFL}}\frac{\Delta x}{c_{\max}},
\end{equation}
где $c_{\max}$~--- максимальная характерная скорость в среде, а $C_{\mathrm{CFL}}$~--- безразмерная константа, зависящая от размерности и вида схемы. Физически это означает, что за один шаг времени информация не должна проходить расстояние, превышающее допустимое для корректного сеточного обмена между узлами. Если это условие нарушается, дискретное решение становится неустойчивым и начинает накапливать нефизические осцилляции.

Для задач термоупругости ограничение по шагу времени может оказаться особенно жёстким, поскольку в системе сочетаются быстрые упругие волны и сравнительно медленные, но чувствительные к сеточной дискретизации тепловые процессы. В сильно неоднородной среде шаг $\Delta t$ приходится выбирать исходя из самых жёстких локальных условий, что увеличивает вычислительные затраты. Кроме того, FDTD подвержен численной дисперсии: если длина волны представлена слишком малым числом узлов, фазовая скорость сеточной волны начинает зависеть от частоты, чего в исходной модели может не быть. Поэтому для высокочастотных режимов и резких контрастов свойств требуются тонкие сетки.

Ещё одно ограничение FDTD связано с геометрией. Для сложных криволинейных границ, разломов, включений и неструктурированных областей регулярная декартова сетка становится неудобной. Использование специальных масок, поглощающих слоёв и модифицированных разностных операторов частично решает проблему, но усложняет схему и может снижать точность вблизи границ.

\subsection*{Спектральные методы}

Спектральные методы основаны на представлении искомого решения в виде разложения по глобальному базису функций. Это могут быть тригонометрические функции, полиномы Чебышёва, Лежандра или другие ортогональные системы. Если, например, решение $u(x,t)$ представлено как
\begin{equation}
u(x,t)=\sum_{m=0}^{M} a_m(t)\phi_m(x),
\end{equation}
то исходная PDE переводится в систему уравнений для коэффициентов $a_m(t)$. Главная идея метода состоит в том, что дифференцирование и интегрирование удобнее выполнять не в физическом пространстве узлов, а в пространстве коэффициентов разложения либо в тесно связанной спектральной области.

Сильная сторона спектральных методов заключается в очень высокой точности для гладких решений. Если коэффициенты среды меняются плавно и сама геометрия области не слишком сложна, ошибка может уменьшаться быстрее, чем у разностных или конечно-элементных методов той же размерности дискретизации. Это делает спектральные подходы особенно привлекательными для параметрических исследований и задач, где важно точно воспроизводить фазовые характеристики волн.

Однако для термоупругих волн в реальных геологических средах спектральные методы сталкиваются с серьёзными ограничениями. Во-первых, резкие скачки коэффициентов на границах слоёв плохо согласуются с глобальными гладкими базисами и могут вызывать осцилляции Гиббса. Во-вторых, сложная геометрия области, разломы, линзы, трещиноватые зоны и криволинейные интерфейсы затрудняют построение эффективного спектрального представления. В-третьих, анизотропия и неоднородность порождают операторы с переменными коэффициентами, для которых реализация спектрального метода становится более громоздкой.

Таким образом, спектральные методы особенно хороши в тех случаях, когда среда достаточно гладкая, область допускает регулярную параметризацию, а основная цель состоит в достижении высокой точности при сравнительно умеренном числе степеней свободы. Но в геологических постановках с сильной структурной сложностью их применимость ограничена.

\subsection*{Метод конечных элементов (FEM)}

Метод конечных элементов строится на другой математической идее: переходе от сильной формы PDE к слабой форме. Смысл этого перехода состоит в том, что искомое решение не обязано обладать классическими производными высокого порядка во всех точках; достаточно, чтобы оно удовлетворяло интегральной форме уравнения на подходящем функциональном пространстве. Для уравнения движения это означает умножение на тестовую функцию $v_i$, интегрирование по области $\Omega$ и использование формулы интегрирования по частям. В результате производные частично переносятся с искомого решения на тестовые функции, а граничные условия на напряжения естественным образом входят в слабую постановку.

Далее область разбивается на конечные элементы, а решение аппроксимируется в каждом элементе локальными базисными функциями:
\begin{equation}
u_h(x,t)=\sum_{a=1}^{N} U_a(t)\,N_a(x), \qquad
T_h(x,t)=\sum_{a=1}^{N} \Theta_a(t)\,N_a(x),
\end{equation}
где $N_a(x)$~--- функции формы, $U_a(t)$ и $\Theta_a(t)$~--- неизвестные узловые коэффициенты. После подстановки в слабую форму задача сводится к системе обыкновенных дифференциальных уравнений по времени для вектора степеней свободы.

Главное преимущество FEM в задачах геофизики состоит в его геометрической гибкости. Метод естественно работает на неструктурированных сетках, допускает сложные границы, локальное сгущение сетки, разрывные и кусочно-заданные коэффициенты, а также удобен для включения анизотропии. Это делает FEM особенно ценным для моделирования термоупругих волн в средах со сложной стратиграфией, разломами и неоднородными включениями.

Кроме того, конечно-элементный подход хорошо подходит для многополевых задач, поскольку механическое и тепловое поля можно аппроксимировать в единой вариационной рамке. Это важно для термоупругой постановки, где уравнения связаны и должны решаться согласованно. Однако высокая универсальность метода достигается ценой большей алгоритмической сложности. Построение сетки, выбор типа элементов, организация интегрирования по элементам, сборка глобальных матриц и решение больших алгебраических систем требуют значительных вычислительных ресурсов. Для нестационарных волновых задач, особенно в трёхмерной постановке, объём вычислений становится очень велик.

\subsection*{Сравнение методов и их ограничения}

Сравнивая FDTD, спектральные методы и FEM, можно отметить, что каждый из них отражает определённую стратегию работы с PDE. FDTD делает ставку на локальную дискретизацию и простоту явной эволюции во времени. Спектральный метод использует глобальное представление решения и достигает высокой точности для гладких задач. FEM переходит к слабой форме и за счёт этого получает максимальную гибкость в отношении геометрии и неоднородности среды.

Для задач термоупругих волн в геологических средах классические методы остаются фундаментальными. Они физически интерпретируемы, математически обоснованы и позволяют получать эталонные решения. Однако именно здесь особенно ярко проявляются их ограничения. Во-первых, высокая размерность трёхмерных задач приводит к колоссальному числу степеней свободы. Во-вторых, сильная неоднородность и многомасштабность требуют очень тонких сеток. В-третьих, анизотропия и сложная геометрия интерфейсов усложняют дискретизацию и повышают вычислительную стоимость. В-четвёртых, при параметрических исследованиях, когда нужно многократно решать задачу для разных наборов свойств среды, прямой численный расчёт становится чрезмерно дорогим.

Именно из этой ситуации возникает естественный интерес к методам искусственного интеллекта. Их задача состоит не в том, чтобы отменить классическую математическую постановку, а в том, чтобы либо ускорить решение, либо научиться эффективно аппроксимировать отображение от параметров среды к полям решения.

\section{Искусственный интеллект в геофизическом моделировании}

Интерес к методам искусственного интеллекта в задачах геофизического моделирования обусловлен прежде всего ограничениями классических вычислительных подходов. Если полная термоупругая задача должна решаться многократно --- например, при сканировании большого пространства параметров, инверсии свойств среды, оптимизационном анализе или формировании больших синтетических выборок, --- даже высокоэффективные классические солверы оказываются слишком затратными. Следовательно, появляется потребность в моделях, которые могли бы либо ускорять прямой расчёт, либо аппроксимировать сам оператор решения.

В геофизике методы машинного обучения применяются в задачах обработки сейсмических данных, распознавания структур, инверсии параметров, аппроксимации решений PDE и построения суррогатных моделей. Для волновых задач особенно важны два класса интеллектуальных подходов. Первый класс ориентирован на приближённое решение уравнений с использованием самой физической постановки как части обучения. Второй класс ориентирован на операторное обучение, то есть на построение отображения между функцией-входом и функцией-решением. В термоупругих задачах оба направления представляют интерес, поскольку исходная модель задаётся системой PDE с переменными коэффициентами и сложными граничными условиями.

\subsection*{Physics-Informed Neural Networks (PINN)}

Идея PINN состоит в том, что нейронная сеть рассматривается как параметрическое приближение искомых полей. Например, можно задать сеть
\begin{equation}
(x,t)\mapsto \big(\hat u_1(x,t),\hat u_2(x,t),\hat u_3(x,t),\hat T(x,t)\big),
\end{equation}
где символом $\hat{\cdot}$ обозначены выходы сети. Производные по координатам и времени вычисляются средствами автоматического дифференцирования, после чего подставляются в уравнения термоупругости. Если обозначить невязки PDE как
\begin{equation}
\mathcal{R}_{\mathrm{mom}}=
\rho \frac{\partial^2 \hat u_i}{\partial t^2}
-\frac{\partial \hat\sigma_{ij}}{\partial x_j},
\qquad
\mathcal{R}_{\mathrm{heat}}=
k\nabla^2\hat T
-C\frac{\partial \hat T}{\partial t}
-\gamma T_0\frac{\partial \hat\varepsilon_{kk}}{\partial t},
\end{equation}
то функция потерь PINN строится как сумма нескольких вкладов:
\begin{equation}
\mathcal{L}
=
\mathcal{L}_{\mathrm{PDE}}
+
\mathcal{L}_{\mathrm{IC}}
+
\mathcal{L}_{\mathrm{BC}}
+
\mathcal{L}_{\mathrm{data}}.
\end{equation}

Здесь $\mathcal{L}_{\mathrm{PDE}}$ штрафует за нарушение уравнений в коллокационных точках внутри области, $\mathcal{L}_{\mathrm{IC}}$ --- за нарушение начальных условий, $\mathcal{L}_{\mathrm{BC}}$ --- за нарушение граничных условий, а $\mathcal{L}_{\mathrm{data}}$ при необходимости учитывает согласование с известными данными.

Математический смысл PINN заключается в том, что поиск решения PDE заменяется задачей оптимизации параметров сети. Сеть не просто подгоняется под примеры, а обучается так, чтобы удовлетворять физическим законам. Для термоупругой задачи это особенно важно, поскольку механические и тепловые поля должны оставаться согласованными. Физически информированная функция потерь позволяет жёстко связать их на этапе обучения.

Преимущества PINN состоят в том, что метод естественно работает с непрерывным представлением решения, относительно просто учитывает сложные граничные условия и потенциально позволяет решать прямые и обратные задачи в одной архитектурной рамке. Для сложной геометрии, если она может быть параметризована, PINN также может быть удобен, поскольку не требует явной сеточной схемы в классическом смысле.

Однако ограничения PINN в задачах волн хорошо известны. Нейросетям трудно воспроизводить высокочастотные осциллирующие решения. Для жёстко связанных многополевых систем, к которым относится термоупругость, баланс между частями функции потерь может оказаться нестабильным. Кроме того, обучение PINN часто требует значительного времени и тщательной настройки весов между невязками уравнений, границ и данных. Поэтому PINN представляют интерес прежде всего как гибкий инструмент для задач умеренной размерности, задач с дефицитом данных и обратных постановок, но не всегда являются самым эффективным способом прямого высокоточного расчёта.

\subsection*{Fourier Neural Operator (FNO)}

FNO относится к классу операторных методов. Его цель состоит не в приближении конкретной функции решения для одной задачи, а в аппроксимации оператора
\begin{equation}
\mathcal{G}: a(x)\mapsto u(x),
\end{equation}
где $a(x)$ обозначает параметры среды, коэффициенты PDE или входное поле, а $u(x)$ --- решение. В термоупругой задаче роль входа могут играть поля $\rho(x)$, $\lambda(x)$, $\mu(x)$, $k(x)$, $C(x)$, $\gamma(x)$, начальные условия и характеристики источника, а роль выхода --- поле смещений и температуры на заданном временном горизонте.

Ключевая идея FNO состоит в том, что операторная свёртка выполняется в спектральной области. На каждом слое сначала применяется преобразование Фурье, затем часть мод умножается на обучаемые комплексные веса, после чего выполняется обратное преобразование. В абстрактной форме один слой можно представить как
\begin{equation}
v_{l+1}(x)=\sigma\Big(Wv_l(x)+\mathcal{F}^{-1}\big(R_l\cdot \mathcal{F}(v_l)\big)(x)\Big),
\end{equation}
где $v_l$ --- скрытое представление на слое $l$, $W$ --- локальное линейное преобразование, $\mathcal{F}$ и $\mathcal{F}^{-1}$ --- прямое и обратное преобразования Фурье, $R_l$ --- обучаемый спектральный множитель, $\sigma$ --- нелинейность.

Преимущество FNO заключается в том, что он учится отображению между функциональными пространствами и способен после обучения быстро выдавать решение для новых конфигураций среды. Для параметрических PDE это чрезвычайно важно: если требуется многократно решать задачу для разных коэффициентов среды, FNO может стать быстрым суррогатным оператором. Кроме того, работа в спектральной области делает метод естественно чувствительным к глобальным структурам решения.

Для термоупругих волн FNO перспективен тогда, когда можно сформировать качественную обучающую выборку из решений классического солвера и когда геометрия области допускает регулярное представление на сетке. Однако метод имеет и ограничения. Во-первых, FNO обычно привязан к сеточному формату данных. Во-вторых, резкие разрывы коэффициентов, сложные локальные интерфейсы и нерегулярные области снижают эффективность стандартной спектральной параметризации. В-третьих, без дополнительных ограничений модель не гарантирует строгого выполнения PDE и граничных условий; она лишь статистически воспроизводит обученный оператор. Поэтому в геофизических приложениях FNO особенно полезен как быстрый суррогат после этапа физически строгого генератора данных.

\subsection*{DeepONet и другие операторные методы}

DeepONet также предназначен для аппроксимации операторов, но делает это в иной архитектурной логике. Основное отображение имеет вид ``функция-вход $\to$ функция-решение''. Обычно входная функция дискретно считывается в нескольких точках и подаётся в ветвь \textit{branch}, а координата точки, в которой требуется значение решения, подаётся в ветвь \textit{trunk}. После этого результат комбинируется, например, скалярным произведением:
\begin{equation}
\hat u(x)=\sum_{q=1}^{Q} b_q(a)\,t_q(x),
\end{equation}
где $b_q(a)$ --- выходы ветви \textit{branch}, зависящие от входной функции $a$, а $t_q(x)$ --- выходы ветви \textit{trunk}, зависящие от координаты $x$. Такая архитектура позволяет рассматривать решение не как фиксированный массив значений, а как значение оператора в произвольной точке области.

От обычной регрессии DeepONet отличается принципиально. В регрессии модель учится отображению из конечномерного вектора в конечномерный вектор. В операторном обучении объектом аппроксимации является отображение между бесконечномерными функциональными пространствами. Для задач математического моделирования это особенно естественно, поскольку коэффициенты среды, источники и начальные условия действительно являются функциями, а не просто наборами чисел.

Для термоупругой задачи DeepONet интересен тем, что позволяет строить модели, которые отображают поля параметров среды в поля решения и потенциально могут быть более гибкими в отношении дискретизации выходного пространства. Однако операторные методы в целом зависят от качества и репрезентативности обучающей выборки. Если пространство возможных геологических конфигураций слишком широко, обучение становится трудным. Кроме того, как и в случае FNO, без специальных добавок физическая строгость решения не гарантируется.

\subsection*{GNN и MeshGraphNet-подобные подходы}

Графовые нейронные сети и MeshGraphNet-подобные методы представляют особый интерес для задач с неструктурированными сетками и сложной геометрией. Если расчётная область разбита на конечные элементы или иные неструктурированные ячейки, её можно представить в виде графа: узлы соответствуют точкам сетки или элементам, а рёбра --- соседству и механико-тепловому взаимодействию. Тогда распространение информации в графовой сети имитирует обмен между соседними фрагментами среды.

Преимущество такого подхода заключается в естественной работе со сложной геометрией, локальными неоднородностями, переменными коэффициентами и разреженными структурами взаимодействий. Для геологических сред, где слои, разломы и включения плохо укладываются в регулярную декартову сетку, это особенно важно. Кроме того, графовое представление хорошо согласуется с конечно-элементной или конечно-объёмной дискретизацией, а значит, открывает путь к гибридным моделям.

Ограничения GNN/MGN связаны с высокой стоимостью обучения на больших графах, трудностями переноса между сетками разной структуры и необходимостью аккуратного кодирования физически значимых признаков узлов и рёбер. Кроме того, без физической регуляризации графовая сеть также остаётся, по сути, аппроксиматором, а не строгим солвером PDE.

\subsection*{Гибридные подходы}

Наиболее перспективным направлением для задач термоупругих волн представляются гибридные подходы, сочетающие классическую физическую постановку с нейросетевыми аппроксимациями. Здесь возможны несколько сценариев. Во-первых, нейросеть может использоваться как суррогатный оператор, который после обучения на данных классического солвера быстро предсказывает решение для новых конфигураций среды. Во-вторых, нейросеть может приближать лишь часть вычислительной цепочки, например эффективный локальный оператор, корректирующий грубосеточное решение. В-третьих, физически информированные сети могут применяться для доуточнения решения, восстановления скрытых коэффициентов среды или ускорения обратной задачи.

Принципиально важно подчеркнуть, что в гибридных схемах физически строгая постановка не исчезает. Напротив, именно классическая система уравнений термоупругости задаёт смысл обучаемого отображения, ограничения на допустимое поведение решения и критерии проверки качества модели. Без этой основы даже самая сложная нейросетевая архитектура остаётся лишь статистическим приближением без гарантии физической корректности.

Следовательно, искусственный интеллект в геофизическом моделировании должен рассматриваться не как замена математической физики, а как расширение вычислительного инструментария. Его основная ценность проявляется там, где классические методы уже обеспечили физическую строгость, но дальнейшая многократная обработка, параметрический анализ или инверсия становятся слишком дорогими.

\section{Выводы по главе 1}

Проведённый анализ показывает, что термоупругие волны в геологических средах представляют собой результат взаимосвязанной динамики механических смещений, напряжений, деформаций и температурных возмущений. В основе математического описания лежит система связанных уравнений: уравнение движения, конститутивное термоупругое соотношение и уравнение теплопереноса с термомеханической связью. Именно эта система отражает физическую сущность процесса: механическое деформирование порождает температурный отклик, а температурное поле, в свою очередь, изменяет напряжённое состояние среды.

Показано, что геологическая среда является сложным параметрическим объектом, а её свойства непосредственно входят в коэффициенты математической модели. Плотность, параметры Ламе, теплопроводность, теплоёмкость, коэффициент температурного расширения, пористость, слоистость и анизотропия не являются внешними характеристиками по отношению к задаче, а формируют саму структуру решения. Контрасты свойств между слоями вызывают отражение, преломление, рассеяние и локализацию энергии. Анизотропия приводит к направленной зависимости скоростей и мод волн. Следовательно, адекватное моделирование термоупругих волн невозможно без учёта реальной структуры геологической среды.

Классические численные методы --- FDTD, спектральные методы и метод конечных элементов --- остаются основой вычислительного решения термоупругих задач. Они обеспечивают физически интерпретируемое и математически обоснованное построение решения. Вместе с тем их применение в реальных геофизических постановках связано с высокими вычислительными затратами, особенно в трёхмерных многомасштабных средах со сложной геометрией и переменными коэффициентами. Именно это обстоятельство делает актуальными методы искусственного интеллекта.

Рассмотренные AI-подходы --- PINN, FNO, DeepONet, GNN/MGN и гибридные схемы --- представляют разные способы аппроксимации решений PDE и операторов решения. Их перспектива заключается в ускорении расчётов, построении суррогатных моделей, решении обратных задач и работе с большим числом параметрических сценариев. Однако эффективность таких методов непосредственно зависит от корректной физической постановки, качества синтетических или экспериментальных данных и способности модели сохранять ключевые законы процесса. Следовательно, наиболее обоснованным направлением является не противопоставление классических и интеллектуальных методов, а их согласованное сочетание.

Тем самым первая глава формирует теоретическую основу для перехода к следующему этапу исследования. Далее логично перейти к конкретной постановке рассматриваемой задачи, выбору геометрии и параметров среды, формулировке начально-краевых условий, построению вычислительной схемы и обоснованию выбранного интеллектуального или гибридного подхода для моделирования термоупругих волн в заданной геологической конфигурации.